\documentclass[12pt,a4paper]{article}

\usepackage[T2A]{fontenc}
\usepackage[utf8]{inputenc}
\usepackage[russian]{babel}
\usepackage{amsmath,amssymb}
\usepackage[margin=2.5cm]{geometry}
\usepackage{enumitem}

\begin{document}

\begin{center}
    {\Large Теоретическое домашнее задание №1}\\[8pt]
\end{center}

\vspace{2em}

\noindent
\textbf{Задача 1.} В аналитическом виде найдите оценки параметра $w$, минимизируя
следующие функции потерь:

\begin{enumerate}[itemsep=6pt, topsep=8pt]
    \item Линейная регрессия без штрафа:
          $L(w) = \frac{1}{n}\sum_{i=1}^{n}(y_i - wx_i)^2$;

    \item Ridge-регрессия:
          $L(w) = \frac{1}{n}\sum_{i=1}^{n}(y_i - wx_i)^2 + \lambda w^2$;

    \item LASSO-регрессия:
          $L(w) = \frac{1}{n}\sum_{i=1}^{n}(y_i - wx_i)^2 + \lambda|w|$;

    \item Пусть решения этих задач равны $\hat w$, $\hat w_R$ и $\hat w_L$ соответственно.
          Найдите пределы
          \[
              \lim_{\lambda \to 0} \hat w_R, \qquad
              \lim_{\lambda \to \infty} \hat w_R, \qquad
              \lim_{\lambda \to 0} \hat w_L, \qquad
              \lim_{\lambda \to \infty} \hat w_L .
          \]

    \item Как можно проинтерпретировать гиперпараметр $\lambda$?
\end{enumerate}

\textbf{Подсказка:} в случае Lasso-регрессии придётся разобраться с модулем. 

\vspace{2em}
\noindent
\textbf{Задача 2.}

Рассмотрим функционал среднего риска для квадратичной функции потерь:

\[
R(a)=\mathbb{E}_{(x, y)\sim p(x, y)}L(y, a(x))=\int_{\mathbb{X}}\int_{\mathbb{Y}}(y-a(x))^2p(x, y)dx dy
\]

Найдите оптимальный алгоритм $a^*(x)$, который минимизирует данный функционал. То есть необходимо решить задачу:

\[
\int_{\mathbb{X}}\int_{\mathbb{Y}}(y-a(x))^2p(x, y)dxdy \rightarrow \min_a
\]

\vspace{2em}
\noindent
\textbf{Задача 3.} 

Убедитесь, что вы знаете ответы на следующие вопросы:
\begin{itemize}
	\item Что такое гиперпараметр модели и чем он отличается от параметра модели?
	\item Почему коэффициент регуляризации нельзя подбирать по обучающей выборке? Как подобрать оптимальное значение для коэффициента регуляризации? 
	\item Почему накладывать регуляризатор на свободный коэффициент~$w_0$ может быть плохой идеей?
	\item Что такое кросс-валидация, чем она лучше использования отложенной выборки?
	\item Почему категориальные признаки нельзя закодировать натуральными числами? Что такое one-hot encoding?
	\item Для чего нужно масштабировать матрицу объекты-признаки перед обучением моделей машинного обучения?
	\item Почему~$L_1$-регуляризация производит отбор признаков?
	\item Почему MSE чувствительно к выбросам? 
\end{itemize}


\end{document}
